\chapter{Optics}
\label{vol04:ch11}

\epigraph{An optical system is a means for transforming wavefronts,
and the operations it performs are limited only by the imagination
of the designer and the discipline of the manufacturer.}{Eugene
Hecht, \emph{Optics}, 5th ed.\ (Pearson, 2017), paraphrased from
Chapter~5 \cite[Chapter~5]{text:hecht-optics}.}

\chapmeta{Half-life: durable. Archetypes: transport, scaling, failure. Exercise target: 30. Project track: physical build plus instrumented measurement (hybrid). Hazard class: medium (laser pointers Class~2, Class~3R; safety-glass and beam-block protocols required). Reader path default: standard, with sections explicitly tagged. Named cases: Hubble Space Telescope primary mirror (1990).}

\noindent
Optics is electromagnetism at wavelengths the human eye can see and at
nearby wavelengths the engineer can manipulate with the same
instruments. The discipline splits into geometric optics (rays
travelling in straight lines, refracting through lenses, reflecting
from mirrors) and wave optics (interference, diffraction, coherence,
polarisation). Both descriptions are correct at the wavelengths and
geometries the engineer cares about; geometric optics is faster when
the relevant geometry is much larger than $\lambda$, wave optics
takes over when geometry approaches $\lambda$. The chapter develops
both, ending with the canonical optical failure mode of the late
twentieth century.

\section{Geometric optics: lenses and mirrors}\pathtag{core}

A \engterm{ray} is the trajectory along which optical energy
propagates in the limit that the wavelength is much smaller than every
relevant geometric scale. Rays bend at refractive interfaces by
Snell's law (Chapter 10), $n_{1}\sin\theta_{1} = n_{2}\sin\theta_{2}$,
reflect at mirrors by the equal-angle rule,
and propagate in straight lines through homogeneous media. Figure
\ref{fig:vol04ch11:snell-refraction} shows the single-interface case
that every lens and prism is built from. Geometric
optics is the working description for camera lenses, telescope
mirrors, microscope objectives, and projection optics; it fails when
diffraction or interference becomes visible, typically when an
aperture's diameter approaches a few wavelengths.

% Figure: a single ray refracting at a plane interface between two media.
% Incident in medium n1 (above), refracted into medium n2 (below), with the
% normal, the incident angle theta1 and the refracted angle theta2 marked.
% Drawn with n2 > n1 so the ray bends toward the normal.
\begin{figure}[htbp]
\centering
\begin{tikzpicture}[>=Stealth, scale=1.0]
% interface (horizontal) and normal (vertical, dashed)
\fill[engblue!6] (-4,0) rectangle (4,-3.2);
\fill[englime!8] (-4,0) rectangle (4,3.2);
\draw[thick] (-4,0) -- (4,0) node[right, font=\small] {interface};
\draw[dashed, enggray] (0,-3.2) -- (0,3.2) node[above, font=\scriptsize] {normal};
% incident ray (upper left, coming down to origin) at 45 deg from normal
\draw[engred, very thick, ->] (-2.6,2.6) -- (-0.05,0.05);
\node[engred, font=\small] at (-2.3,1.6) {incident};
% refracted ray (lower right) bent toward normal, ~28 deg from normal
\draw[engorange, very thick, ->] (0.05,-0.05) -- (1.6,-3.0);
\node[engorange, font=\small] at (1.9,-1.8) {refracted};
% reflected ray (upper right), 45 deg
\draw[enggray, thick, ->, dotted] (0.05,0.05) -- (2.6,2.6);
\node[enggray, font=\scriptsize] at (2.4,1.4) {reflected};
% angle arc theta1 (between normal and incident)
\draw[engred] (0,1.1) arc [start angle=90, end angle=135, radius=1.1];
\node[engred, font=\scriptsize] at (-0.55,1.35) {$\theta_1$};
% angle arc theta2 (between downward normal and refracted)
\draw[engorange] (0,-1.3) arc [start angle=-90, end angle=-62, radius=1.3];
\node[engorange, font=\scriptsize] at (0.5,-1.55) {$\theta_2$};
% medium labels
\node[engblue, font=\small, anchor=west] at (-3.9,1.0) {$n_1$ (e.g.\ air)};
\node[engblue, font=\small, anchor=west] at (-3.9,-1.0) {$n_2 > n_1$ (e.g.\ glass)};
\end{tikzpicture}
\caption[Refraction of a ray at a plane interface]{A ray crossing a plane
interface between media of refractive index $n_1$ and $n_2 > n_1$. Snell's
law, $n_1\sin\theta_1 = n_2\sin\theta_2$, with $n_2 > n_1$ gives
$\theta_2 < \theta_1$, so the transmitted ray bends toward the normal. Part
of the energy reflects at the equal angle $\theta_1$ (dotted). Both angles
are measured from the surface normal, not from the surface.}
\label{fig:vol04ch11:snell-refraction}
\end{figure}


A \engterm{thin lens} is an optical element whose thickness is small
compared to its focal length. The thin-lens equation relates object
distance $s_{o}$ (from the lens), image distance $s_{i}$, and focal
length $f$:
\[
\frac{1}{s_{o}} + \frac{1}{s_{i}} = \frac{1}{f}.
\]
Sign conventions: $s_{o}$ positive for real objects on the incoming
side, $s_{i}$ positive for real images on the outgoing side,
$f$ positive for converging lenses, negative for diverging. The
\engterm{lateral magnification} is
\[
M = -\frac{s_{i}}{s_{o}},
\]
with the negative sign capturing the inversion of a real image cast
by a converging lens. Figure \ref{fig:vol04ch11:lens-ray-diagram}
traces the three principal rays that locate the image of a point and
fixes the sign conventions in a single picture. The calculator
\texttt{code/thin\_lens.py} returns image distance, magnification, and
orientation for any object distance and focal length.

% Figure: thin converging lens ray diagram. Object outside f produces a real,
% inverted, magnified image. Three principal rays: parallel-then-through-F,
% through-center-undeviated, through-F-then-parallel.
% Geometry: f = 2 units; object at s_o = 3 (height 1.5); image at s_i = 6,
% magnification M = -s_i/s_o = -2, image height -3.
\begin{figure}[htbp]
\centering
\begin{tikzpicture}[>=Stealth, scale=0.9]
% optical axis
\draw[enggray, thick] (-4.5,0) -- (8.5,0) node[right, font=\small] {axis};
% lens at x=0
\draw[engblue, very thick] (0,-3.4) -- (0,3.4);
\draw[engblue, ->] (0,3.0) -- (0.35,3.2);
\draw[engblue, ->] (0,3.0) -- (-0.35,3.2);
\draw[engblue, ->] (0,-3.0) -- (0.35,-3.2);
\draw[engblue, ->] (0,-3.0) -- (-0.35,-3.2);
\node[engblue, font=\small, anchor=south] at (0,3.4) {lens};
% focal points f=2
\filldraw[engblack] (-2,0) circle [radius=1.3pt] node[above, font=\scriptsize] {$F$};
\filldraw[engblack] (2,0) circle [radius=1.3pt] node[above right, font=\scriptsize] {$F'$};
% object: at x=-3, height 1.5
\draw[engred, very thick, ->] (-3,0) -- (-3,1.5);
\node[engred, font=\scriptsize, anchor=south] at (-3,1.5) {object};
% image: at x=6, height -3 (inverted)
\draw[engorange, very thick, ->] (6,0) -- (6,-3);
\node[engorange, font=\scriptsize, anchor=north] at (6,-3) {image};
% Ray 1: parallel from object tip to lens, then through F'
\draw[englime!70!black, thick, ->] (-3,1.5) -- (0,1.5);
\draw[englime!70!black, thick, ->] (0,1.5) -- (6,-3);
% Ray 2: through lens center undeviated
\draw[engamber, thick, ->] (-3,1.5) -- (6,-3);
% Ray 3: through F then out parallel
\draw[engblue!70!black, thick, ->] (-3,1.5) -- (0,-1.0);
\draw[engblue!70!black, thick, ->] (0,-1.0) -- (6,-1.0);
\draw[engblue!70!black, thick, ->] (6,-1.0) -- (6,-3);
% labels for s_o and s_i
\draw[<->, enggray, font=\scriptsize] (-3,-0.4) -- (0,-0.4) node[midway, below] {$s_o$};
\draw[<->, enggray, font=\scriptsize] (0,0.4) -- (6,0.4) node[midway, above] {$s_i$};
\end{tikzpicture}
\caption[Thin-lens ray diagram, real inverted image]{Image formation by a
thin converging lens with the object outside the front focal point. Three
principal rays locate the image tip: the ray parallel to the axis refracts
through the back focus $F'$; the ray through the lens centre passes
undeviated; the ray through the front focus $F$ emerges parallel to the
axis. With $f = 2$, $s_o = 3$, the thin-lens equation gives $s_i = 6$ and
magnification $M = -s_i/s_o = -2$: the image is real, inverted, and twice
the object height.}
\label{fig:vol04ch11:lens-ray-diagram}
\end{figure}


The lens-maker's equation gives the focal length of a thin lens of
refractive index $n$ with surface radii $R_{1}$ (first surface) and
$R_{2}$ (second surface),
\[
\frac{1}{f} = (n-1)\left(\frac{1}{R_{1}} - \frac{1}{R_{2}}\right),
\]
under the conventions that $R$ is positive when the centre of
curvature is on the outgoing side. A biconvex lens of radii $R_{1} =
+50\,\si{\milli\meter}$, $R_{2} = -50\,\si{\milli\meter}$ in BK7 glass
($n = 1.52$) has $f = 1/[0.52\times(1/50 + 1/50)] = 48\,\si{\milli
\meter}$. The result is the working tool for first-order optical
design; the actual lens will deviate from the thin-lens prediction
by terms in $(s_{o}/R)^{2}$ and higher (Section 11.2 on aberrations).
The refractive indices and Abbe numbers of the common optical glasses,
the raw material of every lens design, are tabulated in
\texttt{data/glass\_indices.csv}.

A \engterm{spherical mirror} obeys the same equation as a thin lens
in absolute value but with a sign convention specific to reflective
geometry: $f = R/2$ for a concave mirror, with $R$ the radius of
curvature. A concave mirror of radius $1\,\si{\meter}$ has focal
length $0.5\,\si{\meter}$ and produces a real, inverted image of a
distant object at its focal plane. A convex mirror has negative
focal length and produces a virtual, erect, demagnified image used in
car side mirrors and in many security applications. Figure
\ref{fig:vol04ch11:mirror-imaging} traces the image of an object beyond
the centre of curvature of a concave mirror; the construction parallels
the lens diagram, with reflection replacing refraction.

% Figure: concave (converging) spherical mirror imaging an object outside C.
% R = 6 (center of curvature at x=-6), f = R/2 = 3 (focus at x=-3).
% Object at s_o = 8, height 1.2; image real inverted at s_i where
% 1/s_o + 1/s_i = 1/f -> 1/s_i = 1/3 - 1/8 = 5/24 -> s_i = 4.8, M = -0.6.
\begin{figure}[htbp]
\centering
\begin{tikzpicture}[>=Stealth, scale=0.85]
% optical axis (mirror vertex at origin, object side to the left)
\draw[enggray, thick] (-9.2,0) -- (1.2,0) node[right, font=\small] {axis};
% concave mirror arc: vertex at (0,0), opening to the left, R=6 center at (-6,0)
\draw[engblue, very thick] (-6,0) ++(-18:6) arc [start angle=-18, end angle=18, radius=6];
\node[engblue, font=\small, anchor=west] at (0.1,1.7) {mirror};
% center of curvature C and focus F
\filldraw[engblack] (-6,0) circle [radius=1.3pt] node[below, font=\scriptsize] {$C$};
\filldraw[engblack] (-3,0) circle [radius=1.3pt] node[below, font=\scriptsize] {$F$};
% object at x=-8, height 1.2
\draw[engred, very thick, ->] (-8,0) -- (-8,1.2);
\node[engred, font=\scriptsize, anchor=south] at (-8,1.2) {object};
% image at x=-4.8, height -0.72 (inverted, reduced)
\draw[engorange, very thick, ->] (-4.8,0) -- (-4.8,-0.72);
\node[engorange, font=\scriptsize, anchor=north] at (-4.8,-0.78) {image};
% Ray 1: parallel from object tip, reflects through F
\draw[englime!70!black, thick, ->] (-8,1.2) -- (0,1.2);
\draw[englime!70!black, thick, ->] (0,1.2) -- (-4.8,-0.72);
% Ray 2: through C, reflects back on itself (hits mirror normal)
\draw[engamber, thick, ->] (-8,1.2) -- (-4.8,-0.72);
\end{tikzpicture}
\caption[Concave mirror image formation]{A concave spherical mirror of
radius $R = 6$ forms a real, inverted image of an object placed beyond the
centre of curvature $C$. The focal length is $f = R/2 = 3$. A ray parallel
to the axis reflects through the focus $F$; a ray through $C$ strikes the
mirror along its normal and returns on itself. With $s_o = 8$ the mirror
equation $1/s_o + 1/s_i = 1/f$ gives $s_i = 4.8$ and $M = -0.6$: the image
is real, inverted, and reduced, the geometry of a shaving mirror viewed
from far away.}
\label{fig:vol04ch11:mirror-imaging}
\end{figure}


The \engterm{numerical aperture} of an optical system is
\[
\text{NA} = n\sin\theta_{\max},
\]
where $\theta_{\max}$ is the half-angle of the cone of light the
system accepts and $n$ is the refractive index of the working medium
(air, typically; oil immersion in high-NA microscopy). The
$f$-number is the inverse working approximation: $f/\# = 1/(2\,\text{
NA})$ at small angles. A camera lens at $f/2.8$ has NA $\approx 0.18$;
a microscope objective at NA $= 1.4$ (oil immersion) has an
$f/\#$ of about $0.36$.

\section{Aberrations}\pathtag{standard}

Real optical systems deviate from the paraxial (thin-lens, small-
angle) ideal because the spherical surfaces, polychromatic light,
and finite apertures of practical optics produce systematic
distortions called \engterm{aberrations}. Five monochromatic
aberrations are catalogued by the Seidel classification, and two
chromatic aberrations complete the list.

\engterm{Spherical aberration} is the failure of a spherical surface
to focus all parallel rays to a single point: rays striking the
surface farther from the axis focus closer to the lens than paraxial
rays. The transverse spherical aberration scales as $r^{3}$ in the
ray's radial position at the lens. The cure is either an aspheric
surface (a paraboloid focuses all parallel rays to a single point
exactly, used for telescope primary mirrors and high-quality camera
lenses) or a combination of glass elements whose spherical
aberrations cancel by sign.

\engterm{Coma} appears for off-axis points: the image of a point
becomes a comet-shaped streak with its tail pointing radially outward.
Coma scales linearly with field angle and quadratically with aperture.
It is corrected by combining lenses of different shapes (the Abbe
sine condition) or by using off-axis paraboloidal mirrors arranged so
that the on-axis ray reaches the focal point at the correct phase.

\engterm{Astigmatism} of an optical system (distinct from the
eye-anatomical astigmatism) is the failure of a system to focus
tangential and sagittal rays to the same point at off-axis field
positions; tangential rays focus closer to the lens than sagittal
rays, producing an elongated image. Cylindrical lenses introduce or
correct astigmatism deliberately.

\engterm{Field curvature} (Petzval curvature) is the failure of a
flat object plane to image to a flat image plane; the image surface
is a curved bowl. The Petzval condition for flatness requires that
the sum of $1/(n_{i}f_{i})$ over all lens elements be zero, achievable
with a careful selection of converging and diverging elements.

\engterm{Distortion} is the failure of magnification to be constant
across the field; barrel distortion squashes the corners inward,
pincushion distortion stretches them outward. Modern digital cameras
often correct distortion in software, but the optical aberration is
present in the raw image.

\engterm{Chromatic aberration} is the wavelength dependence of the
refractive index: shorter wavelengths (blue) focus closer to the
lens, longer (red) farther. Axial chromatic aberration is the
focal-length variation with wavelength; transverse chromatic
aberration is the magnification variation with wavelength. The cure
is an \engterm{achromatic doublet}: a converging crown-glass element
and a diverging flint-glass element bonded together, with the
combination's chromatic aberration cancelled at two wavelengths
(typically Fraunhofer F at $486\,\si{\nano\meter}$ and C at $656\,\si{
\nano\meter}$). An apochromat extends correction to three wavelengths;
the price is more glass elements and tighter tolerances.

\begin{archetype}[Scaling and failure]
Every aberration scales as a power of $r$ (the ray's radial position
on the lens) and a power of $\theta$ (the field angle). Spherical
aberration is $r^{3}$ at all field positions; coma is $r^{2}\theta$;
astigmatism and field curvature are $r\theta^{2}$; distortion is
$\theta^{3}$. The engineer who knows the scaling can predict which
aberrations dominate at given aperture and field-of-view and direct
the design effort to the dominant terms.
\end{archetype}

\begin{mastery}
\textbf{The Strehl ratio and the Mar\'echal criterion.} A single number
grades how close a real optical system comes to its diffraction limit:
the \engterm{Strehl ratio} $S$, the on-axis peak intensity of the
actual point-spread function divided by the peak of the aberration-free
Airy pattern. For small wavefront errors the extended Mar\'echal
approximation gives $S \approx \exp(-(2\pi\sigma/\lambda)^{2})$, where
$\sigma$ is the root-mean-square wavefront error. The conventional
threshold for ``diffraction-limited'' performance is $S \ge 0.80$, which
the formula reaches at $\sigma = \lambda/14$, the Mar\'echal criterion.
A system at the classical quarter-wave peak-to-valley tolerance sits at
roughly $S = 0.80$ for smooth low-order error. The Hubble primary mirror
of the failure section had a wavefront error near $\lambda/2$ RMS, which
drives $S$ below $0.2$: the formula turns the polishing error directly
into the lost contrast that the first images showed.
\end{mastery}

\section{Wave optics: interference and diffraction}\pathtag{core}

When the wavelength approaches the scale of any obstacle or aperture
in the system, the geometric-optics picture must be extended to
include wave effects. Two phenomena anchor the wave-optics chapter:
interference and diffraction.

\engterm{Interference} occurs when two coherent waves superpose. The
combined intensity is
\[
I = I_{1} + I_{2} + 2\sqrt{I_{1}I_{2}}\cos\delta,
\]
where $\delta$ is the phase difference between the two waves. For
equal-amplitude waves, $I = 4I_{0}\cos^{2}(\delta/2)$, with maxima at
$\delta = 0, 2\pi, 4\pi, \ldots$ and zero at $\delta = \pi, 3\pi,
\ldots$ The \engterm{double-slit experiment} of Thomas Young (1801)
produced the canonical demonstration; the fringe spacing on a screen
at distance $L$ from two slits of separation $d$ is
\[
\Delta y = \frac{\lambda L}{d},
\]
for $\lambda \ll d$. The same fringe-spacing formula governs the
working of every interferometer in metrology. Figure
\ref{fig:vol04ch11:young-double-slit} shows the geometry and the
resulting $\cos^{2}$ intensity pattern.

% Figure: Young double-slit geometry. Two slits separated by d at left, a
% screen at distance L on the right, a point P at height y, and the path
% difference d sin(theta) marked. An intensity sketch sits at the screen.
\begin{figure}[htbp]
\centering
\begin{tikzpicture}[>=Stealth, scale=1.0]
% slit plane (vertical bar at x=0 with two gaps)
\draw[thick] (0,-2.6) -- (0,-0.35);
\draw[thick] (0,-0.15) -- (0,0.15);  % barrier between slits
\draw[thick] (0,0.35) -- (0,2.6);
% two slits at y=+0.25 and y=-0.25
\filldraw[engblue] (0,0.25) circle [radius=1.2pt];
\filldraw[engblue] (0,-0.25) circle [radius=1.2pt];
\draw[<->, enggray, font=\scriptsize] (-0.4,0.25) -- (-0.4,-0.25) node[midway, left] {$d$};
% incoming plane wave (arrows from far left)
\foreach \yy in {-2,-1,0,1,2} {\draw[engred, ->] (-2.2,\yy) -- (-0.4,\yy);}
\node[engred, font=\scriptsize, anchor=south] at (-1.5,2.0) {plane wave};
% screen at x=7
\draw[engblack, very thick] (7,-2.8) -- (7,2.8);
\node[font=\small, anchor=south] at (7,2.8) {screen};
% point P at height 1.6 on screen
\filldraw[engorange] (7,1.6) circle [radius=1.3pt] node[right, font=\scriptsize] {$P$};
% rays from each slit to P
\draw[englime!70!black, thick, ->] (0,0.25) -- (7,1.6);
\draw[engamber, thick, ->] (0,-0.25) -- (7,1.6);
% L marker
\draw[<->, enggray, font=\scriptsize] (0,-2.4) -- (7,-2.4) node[midway, below] {$L$};
% y marker on screen
\draw[<->, enggray, font=\scriptsize] (7.6,0) -- (7.6,1.6) node[midway, right] {$y$};
% intensity sketch to the right of screen
\begin{scope}[shift={(8.6,0)}, scale=1.0]
  \draw[domain=-2.6:2.6, samples=120, smooth, engblue, thick]
    plot (\x, {0.8*cos(deg(\x*2.4))*cos(deg(\x*2.4))});
\end{scope}
\node[engblue, font=\scriptsize, anchor=west] at (8.7,2.4) {fringes};
\end{tikzpicture}
\caption[Young's double-slit geometry]{Two slits separated by $d$,
illuminated by a plane wave, produce overlapping spherical waves that
interfere on a screen a distance $L$ away. The path difference to a point
$P$ at height $y$ is $d\sin\theta \approx dy/L$ for $L \gg d$. Bright
fringes fall where the path difference is an integer number of wavelengths,
$d\sin\theta = m\lambda$, giving a fringe spacing $\Delta y = \lambda L/d$.
The intensity curve at right is the $\cos^2$ interference pattern.}
\label{fig:vol04ch11:young-double-slit}
\end{figure}


A \engterm{thin film} produces interference of a different geometry:
the two interfering beams are the reflection from the film's top surface
and the reflection from its bottom surface, separated by the optical
path $2 n_{f} t\cos\theta_{f}$ through a film of index $n_{f}$ and
thickness $t$ (Figure \ref{fig:vol04ch11:thin-film}). A reflection at a
boundary from a lower index to a higher index adds a $\lambda/2$ phase
shift, which the engineer must track when counting orders. The same
arithmetic that produces the colours of a soap film designs an
anti-reflection coating: a quarter-wave layer of index
$\sqrt{n_{\text{out}} n_{\text{sub}}}$ makes the two reflections cancel
at the design wavelength. The routine \texttt{code/thin\_film.py}
computes the path difference and the quarter-wave thickness.

% Figure: thin-film interference. An incident ray splits into a reflection at
% the top surface and a ray that refracts, reflects off the bottom surface,
% and re-emerges. The two parallel reflected rays interfere.
\begin{figure}[htbp]
\centering
\begin{tikzpicture}[>=Stealth, scale=1.0]
% film of thickness t between y=0 (bottom) and y=1.4 (top)
\fill[engblue!10] (-3.5,0) rectangle (4.5,1.4);
\draw[thick] (-3.5,1.4) -- (4.5,1.4);
\draw[thick] (-3.5,0) -- (4.5,0);
\node[font=\small, anchor=west] at (4.5,0.7) {film, index $n_f$};
\node[font=\scriptsize, anchor=east] at (-3.5,1.9) {medium $n_0 < n_f$};
\node[font=\scriptsize, anchor=east] at (-3.5,-0.45) {substrate $n_s$};
% thickness marker
\draw[<->, enggray, font=\scriptsize] (-3.2,0) -- (-3.2,1.4) node[midway, left] {$t$};
% incident ray from upper left to top surface at point A=(-0.7,1.4)
\draw[engred, very thick, ->] (-2.8,2.7) -- (-0.75,1.45);
\node[engred, font=\scriptsize] at (-2.5,2.1) {incident};
% reflected ray 1 at top surface, goes upper right from A
\draw[engorange, very thick, ->] (-0.65,1.45) -- (1.4,3.0);
\node[engorange, font=\scriptsize] at (0.4,2.6) {ray 1};
% refracted ray into film from A=(-0.7,1.4) down to B on bottom=(0.4,0)
\draw[enggray, thick, ->] (-0.7,1.4) -- (0.4,0);
% reflect at B up to C on top=(1.5,1.4)
\draw[enggray, thick, ->] (0.4,0) -- (1.5,1.4);
% emerge at C as ray 2 (parallel to ray 1)
\draw[englime!70!black, very thick, ->] (1.55,1.45) -- (3.6,3.0);
\node[englime!70!black, font=\scriptsize] at (2.55,2.6) {ray 2};
% mark points
\filldraw[engblack] (-0.7,1.4) circle [radius=1.0pt] node[below left, font=\scriptsize] {$A$};
\filldraw[engblack] (0.4,0) circle [radius=1.0pt] node[below, font=\scriptsize] {$B$};
\filldraw[engblack] (1.5,1.4) circle [radius=1.0pt] node[above left, font=\scriptsize] {$C$};
\end{tikzpicture}
\caption[Thin-film interference]{Interference in a thin film of index $n_f$
and thickness $t$. The incident ray partly reflects at the top surface
(ray~1) and partly refracts, reflects at the bottom surface, and re-emerges
parallel to ray~1 (ray~2). The two reflected rays differ in optical path by
$2 n_f t \cos\theta_f$, plus a $\lambda/2$ phase shift at any reflection
from a lower index to a higher index. Constructive and destructive
interference at each wavelength give the colours of soap films, oil slicks,
and anti-reflection coatings.}
\label{fig:vol04ch11:thin-film}
\end{figure}


\engterm{Diffraction} is the spreading of a wave after passing
through an aperture or around an edge. A wave of wavelength $\lambda$
passing through a slit of width $a$ produces a single-slit
diffraction pattern with first dark fringe at angle
\[
\sin\theta = \frac{\lambda}{a},
\]
and a half-width on a screen at distance $L$ of $\lambda L/a$. The
full intensity profile is the $[\sin\beta/\beta]^{2}$ pattern of Figure
\ref{fig:vol04ch11:single-slit-diffraction}, with a central maximum
twice as wide as the side lobes. A periodic array of slits, a
\engterm{diffraction grating} of period $d$, concentrates the
transmitted light into discrete orders at angles
$d\sin\theta_{m} = m\lambda$ (Figure \ref{fig:vol04ch11:grating-orders});
because the order angle grows with wavelength, a grating disperses white
light into a spectrum and is the working element of most spectrometers.

% Figure: single-slit Fraunhofer diffraction intensity, I = I0 (sin(beta)/beta)^2
% with beta = (pi a / lambda) sin(theta). Plotted against u = (a/lambda) sin(theta),
% so beta = pi*u; zeros at u = +-1, +-2, ... First minima at sin(theta)=lambda/a.
\begin{figure}[htbp]
\centering
\begin{tikzpicture}
\begin{axis}[
  width=12cm, height=7cm,
  xlabel={$(a/\lambda)\sin\theta$},
  ylabel={relative intensity $I/I_0$},
  xmin=-3.5, xmax=3.5, ymin=0, ymax=1.08,
  axis lines=middle,
  xtick={-3,-2,-1,0,1,2,3},
  ytick={0,0.5,1.0},
  tick label style={font=\scriptsize},
  label style={font=\small},
  samples=400, domain=-3.5:3.5,
]
% sinc^2: at x=0 value is 1; guard the x=0 singularity with a tiny offset
\addplot[engblue, very thick] {(sin(deg(pi*x))/(pi*x + 1e-9))^2};
% mark first minima
\draw[dashed, enggray] (axis cs:1,0) -- (axis cs:1,0.25);
\draw[dashed, enggray] (axis cs:-1,0) -- (axis cs:-1,0.25);
\node[engred, font=\scriptsize, anchor=south] at (axis cs:1,0.27)
  {$\sin\theta=\lambda/a$};
\end{axis}
\end{tikzpicture}
\caption[Single-slit diffraction pattern]{Fraunhofer intensity from a single
slit of width $a$, $I = I_0\,[\sin\beta/\beta]^2$ with
$\beta = (\pi a/\lambda)\sin\theta$. The central maximum is twice as wide as
the side lobes and carries most of the energy; the first dark fringes sit at
$\sin\theta = \pm\lambda/a$. Narrowing the slit (smaller $a$) widens the
pattern, the inverse relation between aperture size and diffractive spread
that sets every imaging system's resolution limit.}
\label{fig:vol04ch11:single-slit-diffraction}
\end{figure}


% Figure: diffraction grating. Incident plane wave on a grating of period d;
% diffracted orders m = -2..2 at angles set by d sin(theta_m) = m lambda.
% Geometry illustrative, not to scale.
\begin{figure}[htbp]
\centering
\begin{tikzpicture}[>=Stealth, scale=1.0]
% grating plane (vertical) with tick marks for the rulings
\draw[engblue, very thick] (0,-2.3) -- (0,2.3);
\foreach \yy in {-2,-1.5,-1,-0.5,0,0.5,1,1.5,2}
  {\draw[engblue, thick] (-0.12,\yy) -- (0.12,\yy);}
\node[engblue, font=\small, anchor=south] at (0,2.3) {grating, period $d$};
% incident plane wave from left (normal incidence)
\foreach \yy in {-1.5,-0.75,0,0.75,1.5} {\draw[engred, ->] (-2.4,\yy) -- (-0.3,\yy);}
\node[engred, font=\scriptsize, anchor=south] at (-1.6,1.6) {incident};
% diffracted orders from origin: m=0 straight, m=+-1, m=+-2 at increasing angle
% angles chosen for clarity: m1 = 22 deg, m2 = 48 deg
\def\Lr{4.6}
% m=0
\draw[enggray, very thick, ->] (0,0) -- (\Lr,0);
\node[enggray, font=\scriptsize, anchor=west] at (\Lr,0) {$m=0$};
% m=+1 (22 deg up)
\draw[engorange, thick, ->] (0,0) -- (22:\Lr);
\node[engorange, font=\scriptsize, anchor=south west] at (22:\Lr) {$m=+1$};
% m=-1
\draw[engorange, thick, ->] (0,0) -- (-22:\Lr);
\node[engorange, font=\scriptsize, anchor=north west] at (-22:\Lr) {$m=-1$};
% m=+2 (48 deg)
\draw[englime!70!black, thick, ->] (0,0) -- (48:\Lr);
\node[englime!70!black, font=\scriptsize, anchor=south west] at (48:\Lr) {$m=+2$};
% m=-2
\draw[englime!70!black, thick, ->] (0,0) -- (-48:\Lr);
\node[englime!70!black, font=\scriptsize, anchor=north west] at (-48:\Lr) {$m=-2$};
% angle arc for m=+1
\draw[engorange] (1.4,0) arc [start angle=0, end angle=22, radius=1.4];
\node[engorange, font=\scriptsize] at (24:1.7) {$\theta_1$};
\end{tikzpicture}
\caption[Diffraction grating orders]{A diffraction grating of period $d$ at
normal incidence sends light into discrete orders at angles set by the
grating equation $d\sin\theta_m = m\lambda$. The zeroth order ($m=0$) passes
straight through; higher orders spread to larger angles, and the angle of a
given order grows with wavelength, which is why a grating disperses white
light into a spectrum. The number of visible orders is limited by
$|m| \le d/\lambda$.}
\label{fig:vol04ch11:grating-orders}
\end{figure}


A
circular aperture of diameter $D$ produces an Airy disc with first
dark ring at
\[
\sin\theta = 1.22\frac{\lambda}{D}.
\]
The $1.22$ factor is the first zero of the Bessel function $J_{1}$,
the canonical content of two-dimensional aperture diffraction. The
Airy disc is the angular resolution limit of every telescope, camera,
and microscope; a $200\,\si{\milli\meter}$ telescope at $550\,\si{
\nano\meter}$ resolves angular separations of $3.4\,\si{\micro
\radian}$ ($0.7$ arc-second), the diffraction-limited resolution of
the instrument.

The \engterm{Rayleigh criterion} states that two point sources are
just resolved when the Airy disc of one falls on the first dark ring
of the other; this gives the angular resolution
\[
\theta_{\min} = 1.22\frac{\lambda}{D}.
\]
The criterion is conservative; more sophisticated criteria
(Sparrow, Dawes) recognise smaller resolved separations. The
engineer's working number is the Rayleigh value. Figure
\ref{fig:vol04ch11:rayleigh-resolution} shows two point sources at the
limit, where the summed intensity dips about a quarter below the peaks.
The aperture-resolution trade is tabulated across instruments from the
human eye to the Extremely Large Telescope in
\texttt{data/telescope\_aperture\_resolution.csv}, and the diffraction
calculator \texttt{code/diffraction.py} returns the Rayleigh angle and
the grating orders for any aperture and wavelength.

% Figure: two point sources at the Rayleigh resolution limit. Each is an Airy
% profile (approximated here by a sinc^2-like / gaussian-shaped lobe); the
% peak of one falls on the first minimum of the other. The sum dips ~26%
% between the peaks, the classical Rayleigh "just resolved" criterion.
\begin{figure}[htbp]
\centering
\begin{tikzpicture}
\begin{axis}[
  width=12cm, height=6.6cm,
  xlabel={angular position (units of $1.22\lambda/D$)},
  ylabel={relative intensity},
  xmin=-2.6, xmax=2.6, ymin=0, ymax=1.15,
  axis lines=middle,
  xtick={-2,-1,0,1,2},
  ytick={0,0.5,1.0},
  tick label style={font=\scriptsize},
  label style={font=\small},
  samples=400, domain=-2.6:2.6,
  legend style={font=\scriptsize, at={(0.98,0.98)}, anchor=north east, draw=none},
]
% Airy-like lobe modeled as (2 J1(x)/x)^2 surrogate via (sin/.)^2 scaled.
% Place sources at -0.5 and +0.5 so peaks sit on each other's first zero (zero
% of the surrogate at distance 1.0). Use (sin(pi u)/(pi u))^2 with u offset.
\addplot[engblue, thick, dashed] {(sin(deg(pi*(x+0.5)))/(pi*(x+0.5)+1e-9))^2};
\addplot[engorange, thick, dashed] {(sin(deg(pi*(x-0.5)))/(pi*(x-0.5)+1e-9))^2};
\addplot[engred, very thick]
  {(sin(deg(pi*(x+0.5)))/(pi*(x+0.5)+1e-9))^2 + (sin(deg(pi*(x-0.5)))/(pi*(x-0.5)+1e-9))^2};
\legend{source A, source B, sum}
\end{axis}
\end{tikzpicture}
\caption[Rayleigh resolution criterion]{Two equal point sources at the
Rayleigh limit: the peak of each diffraction pattern falls on the first dark
ring of the other, a separation of $\theta_{\min} = 1.22\lambda/D$ for a
circular aperture of diameter $D$. The summed profile (solid) shows a
central dip of about a quarter below the peaks, the conventional threshold at
which two sources are called ``just resolved.'' Closer than this, the dip
fills in and the pair reads as a single blur. The lobes are drawn with a
$\sin^2$ surrogate for the Airy profile.}
\label{fig:vol04ch11:rayleigh-resolution}
\end{figure}


\begin{estimation}
\textbf{Question.} A spy satellite orbits at $300\,\si{\kilo\meter}$.
What primary-mirror diameter does it need to read a car licence plate
from orbit, and is that diameter physically plausible?

\textbf{Working.} Reading a licence plate means resolving the gaps
between characters, features about $1\,\si{\centi\meter}$ across at a
range of $R = 300\,\si{\kilo\meter}$. The required angular resolution is
$\theta = (0.01\,\si{\meter})/(3\times10^{5}\,\si{\meter}) = 3.3\times
10^{-8}\,\si{\radian}$. The Rayleigh criterion
$\theta = 1.22\lambda/D$ at a visible wavelength
$\lambda = 550\,\si{\nano\meter}$ inverts to
\[
D = \frac{1.22\lambda}{\theta} = \frac{1.22\times 5.5\times10^{-7}}
{3.3\times10^{-8}} \approx 20\,\si{\meter}.
\]
A $20\,\si{\meter}$ space mirror exceeds anything yet flown (Hubble is
$2.4\,\si{\meter}$, the James Webb Space Telescope $6.5\,\si{\meter}$
and segmented), so single-character licence-plate reading from low
orbit is past the diffraction limit of any plausible single aperture.
Relaxing the target to vehicle identification (features
$\sim 0.5\,\si{\meter}$) drops the required diameter to about
$0.4\,\si{\meter}$, which is comfortably within reach, and explains why
real reconnaissance satellites resolve vehicles and structures but not
text. Atmospheric turbulence, neglected here, only makes the text case
worse. The order-of-magnitude verdict, that diffraction alone forbids
the cinematic licence-plate read, holds under every reasonable
assumption.
\end{estimation}

\section{Coherence and lasers (preview)}\pathtag{standard}

Two waves are \engterm{coherent} if they maintain a fixed phase
relationship over time. Coherence is measured along two axes:
temporal coherence (length over which phase is predictable) and
spatial coherence (transverse extent over which phase is predictable).
A thermal source has short coherence in both senses; a laser has
long coherence in both senses, which is why lasers produce
interference fringes with macroscopic path differences and why
thermal sources do not without a beam-splitter and short paths.

The \engterm{coherence length} $L_{c} = c\tau_{c}$ is the distance
over which the wave remains coherent in time. For a source of
linewidth $\Delta f$,
\[
L_{c} \approx \frac{c}{\Delta f}.
\]
A laboratory helium-neon laser has $\Delta f \approx 1.5\,\si{\giga
\hertz}$ (without frequency stabilisation) and $L_{c} \approx 20\,\si{
\centi\meter}$; a frequency-stabilised iodine-locked HeNe has $\Delta
f \le 10\,\si{\kilo\hertz}$ and $L_{c} \ge 30\,\si{\kilo\meter}$. A
fluorescent lamp emits at multiple atomic transitions with $\Delta f$
of order $10^{12}\,\si{\hertz}$ and $L_{c} \approx 300\,\si{\micro
\meter}$; a thermal incandescent bulb has $L_{c}$ of order microns.

The \engterm{laser} (Light Amplification by Stimulated Emission of
Radiation) produces coherent light through stimulated emission in a
gain medium, with feedback provided by a pair of mirrors forming an
optical cavity. The three operating regimes are continuous-wave (the
laser produces a constant output beam), $Q$-switched (the cavity is
held lossy until a population inversion builds, then opened, producing
a giant pulse of $\sim 10\,\si{\nano\second}$ duration), and
mode-locked (multiple cavity modes are forced to oscillate in phase,
producing femtosecond pulses). Working applications: continuous-wave
diode lasers for fibre communications, $Q$-switched solid-state
lasers for laser-induced breakdown spectroscopy, mode-locked
Ti:sapphire lasers for two-photon microscopy and laser materials
processing. The detailed physics is in Chapter 12, where laser
operation is treated as a quantum-mechanical device problem.

\section{Imaging systems: cameras, microscopes, telescopes}\pathtag{standard}

A \engterm{camera} is an imaging system that forms an image of an
extended object on a detector (film or CCD). The camera's specification
is the focal length, the aperture, and the detector size. Focal length
sets the angle of view: a $50\,\si{\milli\meter}$ lens on a
$36\times 24\,\si{\milli\meter}$ full-frame sensor has horizontal
angle of view $\approx 40^{\circ}$, the standard ``normal'' lens; a
$24\,\si{\milli\meter}$ lens gives $74^{\circ}$ (wide-angle); a
$200\,\si{\milli\meter}$ lens gives $10^{\circ}$ (telephoto). The
aperture (relative aperture $f/\#$) sets the working depth of field
and the diffraction limit; an $f/2$ aperture diffraction-limits at
about $1.4\,\si{\micro\meter}$ Airy-disc radius at $550\,\si{\nano
\meter}$. Modern smartphone cameras have aperture $f/1.6$ and pixel
pitch $\approx 1\,\si{\micro\meter}$; the camera is diffraction-
limited rather than sensor-limited at full aperture, and the trick
of computational photography is to recover the lost resolution by
combining multiple exposures with sub-pixel offsets.

A \engterm{microscope} is an imaging system optimised for high
magnification of small objects. The two stages are the objective
(forms an intermediate image at $\sim 160\,\si{\milli\meter}$ from
the back focal plane) and the eyepiece (forms a magnified virtual
image of the intermediate image). The system's resolution is set by
the objective's NA:
\[
d_{\min} = 0.61\frac{\lambda}{\text{NA}}.
\]
At NA $= 1.4$ (oil immersion) and $\lambda = 500\,\si{\nano\meter}$,
$d_{\min} = 220\,\si{\nano\meter}$. Sub-wavelength microscopy
techniques (STED, PALM/STORM, structured illumination) circumvent the
diffraction limit by exploiting nonlinear or molecular-scale
phenomena; their resolution is no longer set by the diffraction limit
of the imaging system but by the spatial precision of the
illumination or the labelling.

A \engterm{telescope} is an imaging system optimised for high angular
resolution of distant objects. The objective (primary mirror or
objective lens) forms an image at its focal plane; the eyepiece
magnifies. The angular resolution is set by the primary's diameter
$D$ via the Rayleigh criterion. The collecting area scales as $D^{2}$;
the resolving power scales as $D$; both are reasons large telescopes
are large.

The largest ground-based optical telescopes have primary mirrors of
$8$ to $10\,\si{\meter}$ diameter (the VLT, Subaru, Keck I and II);
the next generation (Extremely Large Telescope, expected first light
in the late 2020s) will have a $39\,\si{\meter}$ segmented primary.
Atmospheric seeing limits ground-based telescopes to $\sim 1$
arc-second resolution without adaptive optics; with adaptive optics
(a deformable mirror corrects the atmospheric phase distortion in
real time at $\sim 1\,\si{\kilo\hertz}$), large telescopes can reach
their diffraction-limited resolution at the best sites. Space-based
telescopes (Hubble, JWST) avoid atmospheric distortion entirely; the
trade-off is their cost and the difficulty of repair (in Hubble's
case, four servicing missions; in JWST's case, none possible).

\section{Optical fibres}\pathtag{standard}

An \engterm{optical fibre} is a dielectric waveguide consisting of a
high-index core surrounded by a lower-index cladding, with light
confined to the core by total internal reflection (Figure
\ref{fig:vol04ch11:fiber-tir}). The two
predominant fibre types are single-mode (core diameter $\sim 9\,\si{
\micro\meter}$, supporting only the fundamental $\text{LP}_{01}$
mode at $1.55\,\si{\micro\meter}$) and multimode (core diameter
$50\,$or$\,62.5\,\si{\micro\meter}$, supporting many modes).

% Figure: step-index optical fiber. A ray enters within the acceptance cone,
% refracts at the input face, and zig-zags down the core by total internal
% reflection at the core-cladding boundary.
\begin{figure}[htbp]
\centering
\begin{tikzpicture}[>=Stealth, scale=1.0]
% cladding (outer) and core (inner)
\fill[engblue!8] (0,-1.3) rectangle (9,1.3);
\fill[englime!12] (0,-0.7) rectangle (9,0.7);
\draw[thick] (0,-1.3) -- (9,-1.3);
\draw[thick] (0,1.3) -- (9,1.3);
\draw[enggray, thick] (0,-0.7) -- (9,-0.7);
\draw[enggray, thick] (0,0.7) -- (9,0.7);
% input face
\draw[engblack, very thick] (0,-1.3) -- (0,1.3);
% labels
\node[font=\small, anchor=west] at (9.05,0) {core $n_{\text{core}}$};
\node[font=\small, anchor=west] at (9.05,1.0) {cladding $n_{\text{clad}}$};
% acceptance-cone incident ray entering at center-left
\draw[engred, very thick, ->] (-1.9,0.95) -- (-0.05,0.0);
\node[engred, font=\scriptsize] at (-1.5,0.55) {input ray};
% acceptance cone dashed bounds
\draw[engred, dashed] (-1.9,-0.95) -- (0,0);
\draw[engred] (-1.0,-0.5) arc [start angle=-27, end angle=27, radius=1.12];
\node[engred, font=\scriptsize] at (-1.35,0) {$\theta_a$};
% zig-zag inside the core: bounce points along top/bottom of core (y=+-0.7)
% start at (0,0), go up to (1.4,0.7), down to (2.8,-0.7), up to (4.2,0.7)...
\draw[engorange, very thick, ->]
  (0,0) -- (1.4,0.7) -- (2.8,-0.7) -- (4.2,0.7) -- (5.6,-0.7)
       -- (7.0,0.7) -- (8.4,-0.7) -- (9,-0.4);
% normal and TIR angle at one bounce (at (2.8,-0.7))
\draw[dashed, enggray] (2.8,-0.7) -- (2.8,-1.5);
\draw[dashed, enggray] (2.8,-0.7) -- (2.8,0.1);
\node[engorange, font=\scriptsize, anchor=north] at (2.8,-0.72) {$\theta_c$};
\end{tikzpicture}
\caption[Total internal reflection in a step-index fibre]{A ray entering a
step-index fibre within the acceptance cone of half-angle $\theta_a$
refracts at the input face and then strikes the core-cladding boundary at an
angle exceeding the critical angle $\theta_c$, so it reflects totally and
zig-zags down the core with no loss at each bounce. The numerical aperture
$\text{NA} = \sin\theta_a = \sqrt{n_{\text{core}}^2 - n_{\text{clad}}^2}$
sets the largest input angle the fibre will guide.}
\label{fig:vol04ch11:fiber-tir}
\end{figure}


The fibre's \engterm{numerical aperture} is determined by the
refractive-index step:
\[
\text{NA} = \sqrt{n_{\text{core}}^{2} - n_{\text{clad}}^{2}}.
\]
For a typical telecom fibre with $n_{\text{core}} = 1.467$,
$n_{\text{clad}} = 1.460$, $\text{NA} \approx 0.14$, half-angle
$\sim 8^{\circ}$. The routine \texttt{code/fiber\_na.py} returns the
numerical aperture, the critical angle, and the acceptance angle for any
index pair.

The fibre's attenuation has three components:
\begin{itemize}[leftmargin=*]
  \item Rayleigh scattering, dominant at short wavelengths, $\sim
    1/\lambda^{4}$. At $1.55\,\si{\micro\meter}$ it contributes
    $0.16\,\si{\decibel\per\kilo\meter}$.
  \item Hydroxide-ion absorption peaks at $1.39\,\si{\micro\meter}$
    and $1.24\,\si{\micro\meter}$. Modern fibres are dry (low OH$^{-}$
    contamination) and the peaks are largely suppressed.
  \item Infrared lattice absorption rises above $1.6\,\si{\micro
    \meter}$, setting the long-wavelength limit of silica fibres.
\end{itemize}
The window of minimum attenuation is at $1.55\,\si{\micro\meter}$,
about $0.18\,\si{\decibel\per\kilo\meter}$ in modern single-mode
fibres; this is why the long-haul telecom wavelength is
$1.55\,\si{\micro\meter}$ and why erbium-doped fibre amplifiers
(EDFAs) at the same wavelength are the standard amplifier in the
transmission chain. The measured attenuation against wavelength,
showing the residual hydroxide peak and the $1.55\,\si{\micro\meter}$
minimum, is tabulated in \texttt{data/fiber\_attenuation.csv}.

Dispersion in fibres has three components: material (the
wavelength-dependent refractive index of silica), waveguide (the
mode-dependent phase velocity), and modal (in multimode fibre, the
different propagation times of different modes). Material and
waveguide dispersion cancel at the \engterm{zero-dispersion
wavelength} of $\sim 1.31\,\si{\micro\meter}$ in standard single-
mode fibres; this is why $1.31\,\si{\micro\meter}$ is the second
working telecom window. Dispersion-shifted fibres move the zero-
dispersion wavelength to $1.55\,\si{\micro\meter}$ for combined
low-loss and low-dispersion operation.

\section{Worked examples}\pathtag{standard}

Three examples carry the chapter's formulas through to numbers an
engineer would defend.

\paragraph{A camera lens and its diffraction floor.} A
$50\,\si{\milli\meter}$ lens stopped to $f/16$ has an entrance pupil of
$50/16 = 3.1\,\si{\milli\meter}$. The Airy-disc radius at the sensor is
$1.22\lambda(f/\#) = 1.22\times 550\times10^{-9}\times 16 = 10.7\,\si{
\micro\meter}$. A full-frame sensor with $4\,\si{\micro\meter}$ pixels
spreads each point over a five-pixel blur at $f/16$, so stopping down
past about $f/11$ buys depth of field at the cost of resolution. The
optimum aperture for a given sensor balances the growing diffraction
blur against the shrinking aberration blur, and it sits two to three
stops down from wide open for most lenses.

\paragraph{A grating spectrometer's reach.} A grating ruled at
$600\,\si{\per\milli\meter}$ has period $d = 1/600\,\si{\milli\meter} =
1.67\,\si{\micro\meter}$. At normal incidence the first-order angle for
green light at $550\,\si{\nano\meter}$ is
$\theta_{1} = \arcsin(\lambda/d) = \arcsin(0.330) = 19.3^{\circ}$, and
the highest order is $|m| \le d/\lambda = 3.0$, so three orders reach
each side. The angular dispersion $d\theta/d\lambda = m/(d\cos\theta)$
grows with order, which is why high-resolution spectrographs work in
high orders. The values match \texttt{code/diffraction.py} run on
$d = 1.667\,\si{\micro\meter}$.

\paragraph{A fibre link's power budget.} A single-mode link launches
$0\,\si{\decibel}$-referenced power of $1\,\si{\milli\watt}$
($0\,\text{dBm}$) into $80\,\si{\kilo\meter}$ of fibre at
$0.20\,\si{\decibel\per\kilo\meter}$, a span loss of $16\,\si{
\decibel}$. With two connector pairs at $0.5\,\si{\decibel}$ each and a
$3\,\si{\decibel}$ margin, the received power is
$0 - 16 - 1 - 3 = -20\,\text{dBm}$, comfortably above a typical
receiver sensitivity of $-28\,\text{dBm}$ at $10\,\si{\giga\bit\per
\second}$. No mid-span amplifier is needed at this distance; the budget
shows the headroom that lets the same fibre run to $120\,\si{\kilo
\meter}$ before an EDFA becomes necessary.

\begin{mastery}
\textbf{The Abbe diffraction limit as a spatial-frequency cutoff.} The
resolution formula $d_{\min} = 0.61\lambda/\text{NA}$ is better read as
a statement about spatial frequencies than about disc radii. An object
illuminated by a wave carries information in its spatial Fourier
components; a component of period $\Lambda$ diffracts light at angle
$\sin\theta = \lambda/\Lambda$. The objective collects only the
diffracted orders that fall within its acceptance cone, $\sin\theta \le
\text{NA}$, so it passes only spatial frequencies up to
$1/\Lambda_{\min} = \text{NA}/\lambda$, a hard low-pass filter in the
spatial-frequency domain. Finer detail diffracts past the rim of the
objective and never enters the image. This is Abbe's 1873 insight, and
it explains why oil immersion (raising $n$, hence NA above unity) and
shorter wavelengths are the only classical routes to finer resolution,
and why the super-resolution methods of the imaging section must break
the assumptions (linear response, far-field collection) on which the
cutoff rests.
\end{mastery}

\section{Failure: aberrations that ruined a mission}\pathtag{core}

The Hubble Space Telescope launched on 24 April 1990 aboard Space
Shuttle Discovery (STS-31). Within weeks of the first imaging tests,
NASA confirmed that the primary mirror had been figured to the wrong
shape: the outer edge was approximately $2\,\si{\micro\meter}$ too
flat for the design parabola, an error large enough to defeat
diffraction-limited imaging across the telescope's bandpass
\cite{acc:nasa-hubble-optical-systems-1990}.

The NASA Hubble Space Telescope Optical Systems Failure Report,
prepared by the Optical Systems Board of Investigation under Lew
Allen Jr.\ and released in November 1990, traced the error to a
misassembly of the reflective null corrector (RNC), the optical test
instrument used by Perkin-Elmer Corporation to measure the primary
mirror's figure during polishing. A field lens within the RNC was
mispositioned along the optical axis by approximately $1.3\,\si{\milli
\meter}$, the result of a misinterpretation of a measurement target
on a metering rod during the RNC's assembly: the end cap of the rod
and the target feature were not properly distinguished, and the
spacing was set against the wrong feature. The error in the test
instrument propagated into the mirror's polished figure with the
authority of an authoritative test result. Two auxiliary optical
tests, a refractive null corrector and an inverse null test, had
indicated discrepancies during fabrication; both were set aside on
the assumption that the RNC was more authoritative
\cite{acc:nasa-hubble-optical-systems-1990}.

The on-orbit aberration was spherical aberration of $\sim \lambda/2$
peak-to-valley at $550\,\si{\nano\meter}$, severe enough to spread
the Airy disc into a halo of diameter $\sim 1$ arc-second and to
reduce the Strehl ratio (the on-axis peak intensity divided by the
diffraction-limited peak) from a target of $\sim 0.85$ to
approximately $0.15$. Optical correction was delivered by COSTAR
(Corrective Optics Space Telescope Axial Replacement), installed
during the STS-61 servicing mission in December 1993. COSTAR was a
set of small mirrors placed in the optical paths of the
non-camera instruments, figured to cancel the primary's
aberration. The Wide Field/Planetary Camera 2, installed at the same
servicing mission, had its corrective optics integrated into the
instrument itself, the modern approach.

The failure was not the optician's; the Perkin-Elmer team polished
the mirror to the figure their measurement instrument said the
design required. The failure was the absence of process barriers
that would have caused the measurement instrument itself to be
verified, that would have required the auxiliary tests' anomalies to
be formally resolved before fabrication proceeded, and that would
have given NASA's quality-assurance representation effective
oversight of the test apparatus rather than only of the test results.
The Allen Commission's recommendations centred on these missing
barriers: independent verification of test instruments, formal
anomaly-resolution procedures that cannot be informally overridden,
and engaged customer oversight of the test apparatus
\cite{acc:nasa-hubble-optical-systems-1990}.

The Hubble case is the canonical metrology-driven failure in
twentieth-century optical engineering. The lesson by scale of this
chapter is twofold. First, an instrument whose measurement chain is
not independently verified is not a reference; it is a measurement.
Second, when two tests disagree, ``the more authoritative one'' is a
process answer rather than a physical answer; the engineer must
resolve the disagreement before fabricating to either reading.

\begin{principle}[An unverified test instrument is a measurement]
A test instrument that is itself the only authority on its own
measurement chain is not a reference; it is a measurement subject to
the same uncertainty as the artifact it measures. The discipline of
the optical-fabrication shop, and the metrology-traceability
discipline of Volume I, requires that every test instrument's
measurement chain be independently verified before the instrument is
trusted to certify another artifact's figure.
\end{principle}

\section{Project}\pathtag{core}

\begin{project}[Build a Michelson interferometer and measure a small
displacement]
Build a Michelson interferometer on an optical bench and use it to
measure a small displacement at the wavelength precision the
instrument is capable of.

\textbf{Materials.}
A laboratory laser (a $5\,\si{\milli\watt}$, Class~2 or 3R helium-neon
or stabilised diode laser at $632.8\,\si{\nano\meter}$ or
$635\,\si{\nano\meter}$); a $1$-inch beamsplitter cube; two
front-surface mirrors mounted on kinematic bases; a piezoelectric
actuator with a calibrated displacement range of $\sim 10\,\si{
\micro\meter}$ (a stack actuator with $\sim 1\,\si{\volt\per\micro
\meter}$ sensitivity is appropriate); a CCD or smartphone camera as
the imaging detector; an optical breadboard.

\textbf{Procedure.}
\begin{enumerate}
  \item Align the laser, beamsplitter, and two mirrors so that the
    two reflected beams overlap on the detector with a tilt that
    produces fringes of $\sim 1\,\si{\milli\meter}$ spacing.
  \item Adjust the path lengths to within a few millimetres of equal,
    minimising the effect of laser coherence length.
  \item Mount the piezoelectric actuator on the back of one mirror's
    kinematic base so that voltage on the piezo translates the
    mirror along the optical axis.
  \item Apply a triangular voltage waveform to the piezo from $0$ to
    $\pm 10\,\si{\volt}$ and record the fringe motion on the
    detector.
  \item Count the number of fringes that pass a fixed reference line
    per volt of piezo drive. Each fringe corresponds to $\lambda/2$
    of mirror displacement ($316\,\si{\nano\meter}$ for a HeNe).
\end{enumerate}

\textbf{Report.} A one-page report with the setup photograph, the
fringe-counting data, the inferred piezo sensitivity in $\si{\nano
\meter\per\volt}$, the manufacturer's specification, and a discussion
of the discrepancy. Expected: the measured piezo sensitivity should
match the manufacturer's specification to within $\sim 10\,\%$. The
dominant error sources are interferometer-arm mechanical vibration
(reduce by mounting on a vibration-isolated breadboard), thermal
drift of the optical paths (reduce by short measurement intervals),
and detector noise (reduce by averaging multiple sweeps).
\end{project}

\section{Exercises}

\subsection*{Calculation}\pathtag{core}

\begin{exercise}
A thin biconvex lens has surface radii $R_{1} = +75\,\si{\milli\meter}$,
$R_{2} = -50\,\si{\milli\meter}$ in BK7 glass ($n = 1.517$ at
$\lambda = 587\,\si{\nano\meter}$). Compute the focal length.
\paragraph{Solution sketch.} Lens-maker's equation
$1/f = (n-1)(1/R_{1} - 1/R_{2}) = 0.517\,(1/75 - 1/(-50)) = 0.517\,
(0.01333 + 0.02000) = 0.517\times 0.03333 = 0.01723\,\si{\per\milli
\meter}$, so $f = 58.0\,\si{\milli\meter}$. The two convex surfaces add
their curvatures because $R_{2}$ is negative. Matches
\texttt{code/thin\_lens.py}.
\end{exercise}

\begin{exercise}
An object is placed $30\,\si{\centi\meter}$ in front of a converging
thin lens of focal length $20\,\si{\centi\meter}$. Compute the image
distance, the lateral magnification, and whether the image is real
or virtual, erect or inverted.
\paragraph{Solution sketch.} Thin-lens equation
$1/s_{i} = 1/f - 1/s_{o} = 1/20 - 1/30 = 1/60$, so
$s_{i} = 60\,\si{\centi\meter}$. Magnification
$M = -s_{i}/s_{o} = -60/30 = -2$. The positive $s_{i}$ marks a real
image; the negative $M$ marks inversion; $|M| = 2$ marks twofold
enlargement. This is the case drawn in Figure
\ref{fig:vol04ch11:lens-ray-diagram}.
\end{exercise}

\begin{exercise}
A $200\,\si{\milli\meter}$-diameter telescope objective operates at
$550\,\si{\nano\meter}$. Compute the angular resolution (Rayleigh
criterion) in arc-seconds and the minimum separation it can resolve
on the lunar surface ($384\,000\,\si{\kilo\meter}$ distant).
\paragraph{Solution sketch.} Rayleigh angle
$\theta = 1.22\lambda/D = 1.22\times 550\times10^{-9}/0.200 = 3.36\times
10^{-6}\,\si{\radian}$. Converting, $3.36\times10^{-6}\times
(180/\pi)\times 3600 = 0.69$ arc-second. The resolved separation on the
Moon is $\theta R = 3.36\times10^{-6}\times 3.84\times10^{8} = 1.3\,
\si{\kilo\meter}$. A backyard $200\,\si{\milli\meter}$ telescope, if
seeing allowed, could just separate two features $1.3\,\si{\kilo\meter}$
apart on the Moon. Matches \texttt{code/diffraction.py}.
\end{exercise}

\begin{exercise}
A Young's-double-slit experiment uses two slits separated by
$0.5\,\si{\milli\meter}$ at $\lambda = 633\,\si{\nano\meter}$ on a
screen $2\,\si{\meter}$ away. Compute the fringe spacing.
\paragraph{Solution sketch.} Fringe spacing
$\Delta y = \lambda L/d = (633\times10^{-9})(2)/(0.5\times10^{-3}) =
2.53\times10^{-3}\,\si{\meter} = 2.5\,\si{\milli\meter}$. The fringes
are millimetre-scale and easily seen, which is why the double slit is
the classroom demonstration of the wave nature of light. Halving the
slit separation would double the spacing.
\end{exercise}

\begin{exercise}
A microscope objective is rated NA = $0.65$ at the standard $4.7\,\si{
\milli\meter}$ working distance in air. Compute the resolution at
$\lambda = 500\,\si{\nano\meter}$ and the half-angle of the
illumination cone.
\paragraph{Solution sketch.} Resolution
$d_{\min} = 0.61\lambda/\text{NA} = 0.61\times 500\times10^{-9}/0.65 =
4.7\times10^{-7}\,\si{\meter} = 470\,\si{\nano\meter}$. Half-angle from
$\text{NA} = n\sin\theta_{\max}$ in air ($n = 1$):
$\theta_{\max} = \arcsin(0.65) = 40.5^{\circ}$. A dry objective cannot
exceed $\text{NA} = 1$; reaching $470\,\si{\nano\meter}$ resolution
needs no immersion, but finer work does.
\end{exercise}

\begin{exercise}
A single-mode optical fibre has $n_{\text{core}} = 1.4500$,
$n_{\text{clad}} = 1.4467$. Compute the numerical aperture, the
critical angle (from inside the core), and the maximum acceptance
angle in air at the fibre input face.
\paragraph{Solution sketch.}
$\text{NA} = \sqrt{n_{\text{core}}^{2} - n_{\text{clad}}^{2}} =
\sqrt{1.4500^{2} - 1.4467^{2}} = \sqrt{2.1025 - 2.0929} =
\sqrt{0.00957} = 0.0978$. Critical angle (inside core)
$\theta_{c} = \arcsin(n_{\text{clad}}/n_{\text{core}}) =
\arcsin(0.99772) = 86.1^{\circ}$, so rays must strike the boundary
within $3.9^{\circ}$ of grazing to be guided. Acceptance half-angle in
air $\theta_{a} = \arcsin(\text{NA}) = \arcsin(0.0978) = 5.6^{\circ}$.
Matches \texttt{code/fiber\_na.py}.
\end{exercise}

\begin{exercise}
A laser of wavelength $785\,\si{\nano\meter}$ has a linewidth of
$10\,\si{\mega\hertz}$. Compute the coherence length and the
maximum interferometer-arm imbalance for which fringes can still
form.
\paragraph{Solution sketch.} Coherence length
$L_{c} \approx c/\Delta f = (3\times10^{8})/(10\times10^{6}) =
30\,\si{\meter}$. Fringe contrast falls off as the path imbalance
approaches $L_{c}$; usable fringes persist while the round-trip
imbalance is well within $L_{c}$, so the arm-length difference (half the
round trip) can be up to about $15\,\si{\meter}$. The linewidth, not the
wavelength, sets how far apart the arms may be, which is why narrow-line
lasers are the metrology choice.
\end{exercise}

\subsection*{Derivation}\pathtag{standard}

\begin{exercise}
Derive the thin-lens equation $1/s_{o} + 1/s_{i} = 1/f$ by applying
the small-angle refraction equation $n_{1}/s_{o} + n_{2}/s_{i} =
(n_{2}-n_{1})/R$ at both surfaces of a thin lens in air. State
precisely where the thin-lens approximation is used.
\paragraph{Solution sketch.} Apply the single-surface refraction
relation at surface~1 (air to glass, radius $R_{1}$):
$1/s_{o} + n/s_{i1} = (n-1)/R_{1}$, where $s_{i1}$ is the intermediate
image distance. The intermediate image serves as the object for
surface~2 (glass to air, radius $R_{2}$); for a lens of zero thickness
the object distance there is $-s_{i1}$ (a virtual object on the far
side), giving $-n/s_{i1} + 1/s_{i} = (1-n)/R_{2}$. Adding the two
equations cancels the $s_{i1}$ terms and yields
$1/s_{o} + 1/s_{i} = (n-1)(1/R_{1} - 1/R_{2}) = 1/f$. The thin-lens
approximation enters at exactly one step: setting the object distance
for surface~2 equal to $-s_{i1}$, which assumes the lens thickness is
negligible compared with the image distances.
\end{exercise}

\begin{exercise}
Derive the Rayleigh resolution criterion $\theta_{\min} = 1.22\lambda
/D$ for a circular aperture by treating the aperture function as
unity inside a disc and zero outside, computing its Fraunhofer
diffraction pattern, and identifying the first zero of the resulting
Bessel function.
\paragraph{Solution sketch.} The Fraunhofer pattern is the
two-dimensional Fourier transform of the aperture function. For a
uniform disc of radius $a$, the transform is the circularly symmetric
integral $\int_{0}^{a} J_{0}(kr\sin\theta)\,r\,dr$, which evaluates to
$\propto J_{1}(ka\sin\theta)/(ka\sin\theta)$ using the identity
$\int_{0}^{x} J_{0}(u)u\,du = x J_{1}(x)$. The intensity is the square
of this amplitude, the Airy pattern. The first zero is the first zero of
$J_{1}$, at argument $3.832$. Setting
$ka\sin\theta = (2\pi/\lambda)(D/2)\sin\theta = 3.832$ and solving gives
$\sin\theta = 3.832\lambda/(\pi D) = 1.22\lambda/D$, since
$3.832/\pi = 1.22$. The Rayleigh criterion places the second source's
peak at this first zero.
\end{exercise}

\begin{exercise}
Derive the lens-maker's equation $1/f = (n-1)(1/R_{1} - 1/R_{2})$
from the small-angle refraction relation, applied twice.
\paragraph{Solution sketch.} Send parallel light ($s_{o} \to \infty$)
into the thin lens; the image forms at the focal point, $s_{i} = f$. The
two-surface derivation of the previous problem reduces, with
$1/s_{o} = 0$, to $1/f = (n-1)(1/R_{1} - 1/R_{2})$ directly. The same
cancellation of the intermediate image distance does the work; the
infinite-object case isolates $f$ without carrying the imaging terms,
and the result holds for finite objects because $f$ is a property of the
lens alone.
\end{exercise}

\begin{exercise}
Derive the fringe-spacing formula $\Delta y = \lambda L/d$ for
Young's-double-slit setup, starting from the path-length difference
between the two slits and a far-field point on the screen.
\paragraph{Solution sketch.} The path difference from the two slits to a
screen point at angle $\theta$ is $\delta = d\sin\theta$. Bright fringes
require $\delta = m\lambda$, so $\sin\theta_{m} = m\lambda/d$. For
$L \gg d$ the small-angle approximation gives
$\sin\theta \approx \tan\theta = y/L$, so the $m$th bright fringe sits
at $y_{m} = m\lambda L/d$. Adjacent fringes differ by $m \to m+1$, so
$\Delta y = y_{m+1} - y_{m} = \lambda L/d$. The spacing is uniform
because the path difference is linear in $y$ in this approximation. See
Figure \ref{fig:vol04ch11:young-double-slit}.
\end{exercise}

\begin{exercise}
Derive the numerical-aperture formula $\text{NA} = \sqrt{n_{\text{
core}}^{2}-n_{\text{clad}}^{2}}$ for a step-index optical fibre, by
considering a ray at the critical angle for total internal reflection
at the core-cladding boundary and tracing it back to the fibre's
input face.
\paragraph{Solution sketch.} A ray guided by total internal reflection
strikes the core-cladding boundary at the critical angle
$\theta_{c}$, where $\sin\theta_{c} = n_{\text{clad}}/n_{\text{core}}$.
The complementary angle to the fibre axis is
$\phi = 90^{\circ} - \theta_{c}$, with
$\sin\phi = \cos\theta_{c} = \sqrt{1 - (n_{\text{clad}}/n_{\text{core}})
^{2}}$. Refraction at the input face (air to core) gives
$\sin\theta_{a} = n_{\text{core}}\sin\phi$ by Snell's law, so
$\sin\theta_{a} = n_{\text{core}}\sqrt{1 - n_{\text{clad}}^{2}/
n_{\text{core}}^{2}} = \sqrt{n_{\text{core}}^{2} - n_{\text{clad}}^{2}}$.
That is the numerical aperture, $\text{NA} = \sin\theta_{a}$. See Figure
\ref{fig:vol04ch11:fiber-tir}.
\end{exercise}

\subsection*{Estimation}\pathtag{core}

\begin{exercise}
Estimate the minimum reliably-detected angular separation of two
stars by the unaided human eye (pupil diameter $\sim 4\,\si{\milli
\meter}$, peak visual wavelength $550\,\si{\nano\meter}$). Compare
the diffraction limit to the empirical visual acuity of about
$1$ arc-minute.
\paragraph{Solution sketch.} Diffraction limit
$\theta = 1.22\lambda/D = 1.22\times 550\times10^{-9}/0.004 =
1.7\times10^{-4}\,\si{\radian}$, which is
$1.7\times10^{-4}\times(180/\pi)\times 60 = 0.58$ arc-minute. The eye's
empirical acuity of about $1$ arc-minute is within a factor of two of
the diffraction limit, so the eye is close to diffraction-limited at the
centre of the field. Retinal cone spacing (about $2\,\si{\micro\meter}$
at the fovea, subtending roughly $0.5$ arc-minute) is matched to the
diffraction limit, the expected outcome of a well-engineered detector.
\end{exercise}

\begin{exercise}
Estimate the temporal coherence length of sunlight falling on the
ground, given that visible sunlight covers a wavelength range from
$\sim 400$ to $\sim 700\,\si{\nano\meter}$. State the implication for
laboratory interferometry without a narrowband filter.
\paragraph{Solution sketch.} The frequency spread of visible sunlight is
$\Delta f = c(1/\lambda_{\min} - 1/\lambda_{\max}) = 3\times10^{8}\,
(1/400\times10^{-9} - 1/700\times10^{-9}) \approx 3.2\times10^{14}\,\si{
\hertz}$. The coherence length is
$L_{c} \approx c/\Delta f \approx 3\times10^{8}/3.2\times10^{14} \approx
0.9\,\si{\micro\meter}$, of order one wavelength. Interferometry with
white light works only when the two arms are matched to within about a
micron, which is why white-light fringes are used to locate the
zero-path-difference point precisely, and why a narrowband filter is
needed to see fringes at any appreciable path imbalance.
\end{exercise}

\begin{exercise}
Estimate the depth of field of a $50\,\si{\milli\meter}$ camera lens
focused at $5\,\si{\meter}$ at $f/4$, with the circle of confusion
set to $30\,\si{\micro\meter}$ (the diagonal divided by $\sim 1500$
for a full-frame sensor).
\paragraph{Solution sketch.} The hyperfocal distance is
$H \approx f^{2}/(Nc) = (0.050)^{2}/(4\times 30\times10^{-6}) =
0.0025/1.2\times10^{-4} = 20.8\,\si{\meter}$. With the lens focused at
$s = 5\,\si{\meter}$, the near and far limits are
$D_{n} = sH/(H+s) = 5\times 20.8/25.8 = 4.0\,\si{\meter}$ and
$D_{f} = sH/(H-s) = 5\times 20.8/15.8 = 6.6\,\si{\meter}$, a total depth
of field of about $2.6\,\si{\meter}$. The depth is asymmetric, extending
farther behind the focus than in front, the familiar one-third-in-front
rule of thumb.
\end{exercise}

\begin{exercise}
Estimate the photon-shot-noise limit of a $14$-bit, $5\,\si{\micro
\meter}$-pixel scientific CMOS sensor with $50\,000$-electron full
well at maximum brightness. State the consequent dynamic-range
limit and the read-noise floor that would be acceptable.
\paragraph{Solution sketch.} At full well $N = 50\,000$ electrons, the
shot noise is $\sqrt{N} = 224$ electrons, so the signal-to-noise ratio
at saturation is $50\,000/224 = 224$. The shot-noise-limited dynamic
range is therefore about $224{:}1$, or $47\,\si{\decibel}$. Mapping the
full well onto $14$ bits gives $50\,000/16384 = 3.05$ electrons per
digital level, so a read-noise floor near $3$ electrons matches the
quantisation and the digitiser is well chosen; a read noise much below
$3$ electrons would be wasted because quantisation dominates, and much
above would throw away the sensor's low-light reach.
\end{exercise}

\subsection*{Simulation}\pathtag{mastery}

\begin{exercise}
Implement a Fresnel-propagation algorithm using the
angular-spectrum method to propagate a Gaussian beam from a $\lambda
/100$ initial spot to a $1\,\si{\meter}$-distant observation plane.
Verify against the closed-form Gaussian-beam propagation and
identify the regime in which the Fresnel approximation begins to
break down.
\paragraph{Solution sketch.} The angular-spectrum method takes the
two-dimensional Fourier transform of the initial field, multiplies by
the free-space transfer function $H(f_{x},f_{y}) = \exp(i k z
\sqrt{1 - (\lambda f_{x})^{2} - (\lambda f_{y})^{2}})$, and
inverse-transforms. Sample the input plane finely enough that the
highest spatial frequency $\text{NA}/\lambda$ is below the Nyquist
limit; for a $\lambda/100$ spot this demands a grid pitch well under
$\lambda$. Compare the propagated waist and on-axis amplitude with the
closed-form Gaussian beam, $w(z) = w_{0}\sqrt{1 + (z/z_{R})^{2}}$ with
$z_{R} = \pi w_{0}^{2}/\lambda$. The Fresnel (paraxial) approximation
expands the square root to first order and breaks down when
$(\lambda f)^{2}$ is no longer small, that is for high-NA beams or very
short propagation where evanescent and wide-angle components matter;
agreement should be excellent for the modest-NA Gaussian and degrade as
$w_{0}$ shrinks toward $\lambda$.
\end{exercise}

\begin{exercise}
Implement a ray-tracer for a thin lens with spherical aberration,
sampling rays uniformly across the aperture. Plot the spot diagram
at the paraxial focus, the marginal focus, and the circle of least
confusion. Identify the empirical optimum focus position relative to
the paraxial focus.
\paragraph{Solution sketch.} For a thin lens with spherical aberration,
a ray entering at radial height $r$ crosses the axis at a longitudinal
position shifted from the paraxial focus by $\Delta z \propto r^{2}$
(transverse aberration $\propto r^{3}$). Sample rays uniformly in $r$,
trace each to the chosen image plane, and plot the intersection points.
At the paraxial focus the marginal rays form a large outer ring around a
bright core; at the marginal focus the core blurs and the marginal rays
tighten. The circle of least confusion, the smallest enclosing spot,
sits about three-quarters of the way from the paraxial focus toward the
marginal focus for third-order spherical aberration; that intermediate
plane is the empirical best focus, and its spot radius is roughly a
quarter of the paraxial-plane ring radius.
\end{exercise}

\subsection*{Design}\pathtag{standard}

\begin{exercise}
Design a $20\times$ achromatic microscope objective at a working
distance of $5\,\si{\milli\meter}$ in air. Specify the focal length,
the aperture, the expected NA, and the back-focal-plane diameter
required for an eyepiece of $f = 10\,\si{\milli\meter}$.
\paragraph{Solution sketch.} For a standard $160\,\si{\milli\meter}$
tube length, a $20\times$ objective has focal length
$f_{\text{obj}} = 160/20 = 8\,\si{\milli\meter}$. A $20\times$ objective
typically carries $\text{NA} \approx 0.40$, so the half-angle is
$\arcsin(0.40) = 23.6^{\circ}$ and the front aperture diameter is about
$2 f_{\text{obj}}\tan(23.6^{\circ}) \approx 7\,\si{\milli\meter}$. The
achromat is a cemented crown-flint doublet correcting axial colour at
the F and C lines. The back-focal-plane (exit-pupil) diameter that fills
a $10\,\si{\milli\meter}$ eyepiece is set by matching exit pupil to
eyepiece field stop, of order $4$ to $5\,\si{\milli\meter}$ for
comfortable eye relief. The working distance of $5\,\si{\milli\meter}$
is short but consistent with an $8\,\si{\milli\meter}$ focal length.
\end{exercise}

\begin{exercise}
Design an optical system to image a $1\,\si{\centi\meter}$-square
object onto a $36\times 24\,\si{\milli\meter}$ camera sensor at unit
magnification. Specify the lens focal length, the working distance,
the aperture, and the expected diffraction limit at $f/8$.
\paragraph{Solution sketch.} Unit magnification ($M = -1$) requires
object and image each at $2f$, so the object and sensor are $4f$ apart.
A $1\,\si{\centi\meter}$ object onto a $36\times24\,\si{\milli\meter}$
sensor at $M = -1$ over-fills the long axis only slightly; choose a
convenient focal length, say $f = 100\,\si{\milli\meter}$, giving object
distance $s_{o} = 200\,\si{\milli\meter}$, image distance
$s_{i} = 200\,\si{\milli\meter}$, and a working distance near
$200\,\si{\milli\meter}$. At $f/8$ the effective aperture at $M = -1$ is
$N_{\text{eff}} = N(1+|M|) = 16$, so the Airy radius is
$1.22\lambda N_{\text{eff}} = 1.22\times 550\times10^{-9}\times 16 =
10.7\,\si{\micro\meter}$. The bellows factor doubling the effective
$f$-number is the detail a unit-magnification macro design must not
omit.
\end{exercise}

\begin{exercise}
Design a single-mode-fibre link for $10\,\si{\giga\bit\per\second}$
data transmission over $80\,\si{\kilo\meter}$ at $1.55\,\si{\micro
\meter}$, using one EDFA amplifier mid-span. Specify the launch
power, the receiver sensitivity, the link budget, the amplifier gain,
and the dispersion compensation required.
\paragraph{Solution sketch.} Span loss at
$0.20\,\si{\decibel\per\kilo\meter}$ over $80\,\si{\kilo\meter}$ is
$16\,\si{\decibel}$; add $\sim 2\,\si{\decibel}$ for connectors and
splices and a $3\,\si{\decibel}$ margin, total $21\,\si{\decibel}$. With
a launch power of $+3\,\text{dBm}$ and a receiver sensitivity of
$-24\,\text{dBm}$ at $10\,\si{\giga\bit\per\second}$, the budget is
$+3 - 21 = -18\,\text{dBm}$ received, $6\,\si{\decibel}$ above
sensitivity, so a mid-span EDFA with about $12\,\si{\decibel}$ gain
restores margin and supports growth. Chromatic dispersion of standard
single-mode fibre is $\sim 17\,\si{\pico\second/(\nano\meter\,\kilo
\meter)}$; over $80\,\si{\kilo\meter}$ with a $0.1\,\si{\nano\meter}$
source linewidth the spread is $\sim 136\,\si{\pico\second}$, a
significant fraction of the $100\,\si{\pico\second}$ bit period, so a
dispersion-compensating fibre module of about $-1360\,\si{\pico\second/
\nano\meter}$ is specified to flatten it.
\end{exercise}

\subsection*{Diagnosis and reverse engineering}\pathtag{standard}

\begin{exercise}
A microscope's spot diagram at a single point of the field of view
shows a triangular three-lobed pattern instead of an Airy disc.
Identify the candidate aberrations (trefoil from a three-point mirror
support, primary trefoil from a manufacturing error) and the
supplementary measurement that distinguishes them.
\paragraph{Solution sketch.} A three-lobed spot is the signature of
trefoil, a $\cos 3\phi$ wavefront error. Mount-induced trefoil from a
three-point support rotates with the optic when the mount is reclocked;
a manufacturing (figure) trefoil is frozen into the glass and stays
fixed in the optic's own frame. The distinguishing measurement is to
rotate the suspect element relative to its mount and watch the lobe
orientation: if the lobes follow the mount, the support is the cause
(relieve or re-space the contacts); if they follow the glass, the figure
is at fault (re-polish or reject). An interferogram or Shack-Hartmann
wavefront map confirms the $\cos 3\phi$ term quantitatively.
\end{exercise}

\begin{exercise}
An astronomical CCD image shows two distinct ring patterns around
bright stars: an inner Airy ring and an outer halo. Identify the
candidate optical defects (defocused secondary mirror, primary-mirror
spherical aberration, dome seeing) and the test that distinguishes
them.
\paragraph{Solution sketch.} A defocused secondary throws a symmetric
ring whose diameter changes sharply with focus and which collapses to a
point at best focus; the test is to step through focus and watch the
halo shrink to nothing. Primary-mirror spherical aberration leaves a
residual halo that cannot be focused away at any setting, the Hubble
signature; the test is that no focus position removes it, and a
star-test or interferogram shows the $r^{4}$ wavefront term. Dome (local
air) seeing produces a halo that varies in time and with thermal
conditions, worst soon after dome opening; the test is to monitor the
pattern over minutes and correlate with temperature, and it improves as
the air settles. Persistence under focus and over time separates the
three.
\end{exercise}

\begin{exercise}
A laboratory measurement of an interferometer's fringe contrast
shows that the contrast drops as the path-length difference exceeds
$10\,\si{\centi\meter}$, even though the source is specified as a
laboratory HeNe laser. Identify the candidate causes (multi-mode
laser operation, vibration washing out the fringe, mechanical drift)
and the test that distinguishes them.
\paragraph{Solution sketch.} A contrast drop that sets in at a definite
path difference is the signature of finite coherence length, which an
unstabilised multi-longitudinal-mode HeNe shows at the
$\sim 20\,\si{\centi\meter}$ scale of its mode beating; the test is to
measure contrast against path difference and look for periodic revivals
at multiples of the cavity round-trip length, the fingerprint of mode
structure. Vibration washes the fringes out at all path differences and
is time-varying, not path-dependent; the test is to isolate the bench or
shorten the integration and watch the contrast recover. Mechanical drift
moves the fringes slowly and blurs only on long exposures; the test is
to shorten the exposure and see the instantaneous contrast return. The
path-difference dependence cleanly fingers coherence length as the
cause.
\end{exercise}

\subsection*{Failure analysis}\pathtag{core}

\begin{exercise}
Summarise the technical mechanism of the Hubble Space Telescope
primary mirror failure (1990) in three sentences, using the registry
entry's working language. State which process barrier was missing
and which auxiliary tests were set aside.
\paragraph{Solution sketch.} A field lens in the reflective null
corrector used to measure the primary mirror was mispositioned by about
$1.3\,\si{\milli\meter}$, the result of setting a spacing against the
wrong feature on a metering rod. The faulty test instrument certified
the mirror as correct while it was being polished about
$2\,\si{\micro\meter}$ too flat at the edge, producing on-orbit
spherical aberration of order $\lambda/2$. The missing process barrier
was independent verification of the test instrument's own measurement
chain; two auxiliary tests, a refractive null corrector and an inverse
null test, had flagged discrepancies and were set aside as less
authoritative than the reflective null corrector
\cite{acc:nasa-hubble-optical-systems-1990}.
\end{exercise}

\begin{exercise}
A $300\,\si{\milli\meter}$-diameter astronomical telescope shows a
visible coma pattern in off-axis stars after assembly. Identify the
candidate failure modes (decentred primary, tilted secondary, focal-
plane misalignment) and the test that would distinguish them.
\paragraph{Solution sketch.} Field-uniform coma that points the same way
everywhere off-axis indicates a misaligned (decentred or tilted) optic,
since a centred system has zero coma on-axis and coma growing linearly
with field. A decentred primary and a tilted secondary both add a
constant coma across the field; they are separated by a star test on
axis, where a centred system shows a symmetric Airy pattern: residual
on-axis coma points to the secondary tilt, while a coma that reverses
sense when the secondary is reclocked points to the secondary. Focal-
plane (detector) misalignment produces defocus or astigmatism varying
across the frame, not uniform coma; the test is to check focus at
several field points. Collimation by the standard star test, adjusting
the secondary until on-axis coma vanishes, isolates and cures the
dominant term.
\end{exercise}

\begin{exercise}
A high-power laser-cutting system fails to maintain its specified
focal-spot size after two years of operation, even though all
maintenance procedures have been followed. Identify the candidate
failure modes (thermal lensing in the focusing lens, contamination
of the output window, beam-quality degradation of the laser source)
and the working diagnostic test for each.
\paragraph{Solution sketch.} Thermal lensing in the focusing optic
shifts the focal plane as the lens heats during a cut; the diagnostic is
to measure spot size against dwell time or duty cycle and look for a
drift that grows with absorbed power, then inspect the lens for
absorbing contamination. Output-window contamination scatters and
absorbs power at the window; the diagnostic is a power-meter reading
before and after the window plus a microscope inspection for deposits or
pitting. Source beam-quality degradation raises the beam parameter
product $M^{2}$, enlarging the focal spot independently of the delivery
optics; the diagnostic is a beam-profiler $M^{2}$ measurement at the
source output compared against the commissioning value. The
power-dependent (thermal), window-localised (contamination), and
source-intrinsic ($M^{2}$) signatures each point to one element.
\end{exercise}

\subsection*{Judgment}\pathtag{standard}

\begin{exercise}
A consulting client must decide between an aspheric singlet and a
multi-element achromat for a $5\,\si{\milli\meter}$-focal-length
collimator at $\lambda = 633\,\si{\nano\meter}$. Identify the trade-
offs in terms of cost, performance, weight, and field of view.
\paragraph{Solution sketch.} At a single design wavelength
($633\,\si{\nano\meter}$) chromatic correction is not needed, so the
aspheric singlet wins on weight, element count, and on-axis spot size,
and a moulded asphere is cheap in volume though costly in small runs.
The achromat carries no advantage at one wavelength but tolerates a
spread of wavelengths and often shows better off-axis behaviour (a
singlet asphere corrects spherical aberration on-axis but leaves coma
and field curvature). The recommendation: for a monochromatic on-axis
collimator the asphere is the right call; choose the achromat only if
the source is broadband, the field is wide, or small-quantity asphere
tooling cost dominates.
\end{exercise}

\begin{exercise}
A research group must choose between conventional optical microscopy
at NA = $1.4$ and super-resolution microscopy at $20\,\si{\nano
\meter}$ resolution for a study of cellular ultrastructure. Identify
the trade-offs in cost, sample preparation, time, and biological
relevance.
\paragraph{Solution sketch.} Conventional confocal at $\text{NA} = 1.4$
resolves about $220\,\si{\nano\meter}$, images live cells with standard
fluorophores, is fast, and is widely available; it cannot see structure
below the diffraction limit. Super-resolution at $20\,\si{\nano\meter}$
(STED, PALM/STORM) reaches the molecular scale but demands special
fluorophores or photoswitching, longer acquisition (especially
single-molecule localisation), higher cost and expertise, and often
fixed rather than live samples, raising the question of whether the
fixed-cell ultrastructure reflects the living state. The recommendation
turns on the science: if the target features are above
$\sim 250\,\si{\nano\meter}$ or live dynamics matter, conventional
microscopy serves; if the question is specifically sub-diffraction
ultrastructure and fixation is acceptable, the super-resolution cost is
justified.
\end{exercise}

\begin{exercise}
A research group must decide whether to deploy a $1\,\si{\meter}$
ground-based telescope with adaptive optics or to wait for time on
the JWST for an exoplanet-imaging programme. Identify the trade-offs
in resolution, sensitivity, schedule, and risk.
\paragraph{Solution sketch.} A $1\,\si{\meter}$ ground telescope with
adaptive optics reaches a diffraction limit of
$\sim 0.14$ arc-second in the near-infrared, is available on the group's
own schedule, and is cheap per night, but it loses sensitivity to sky
background and residual atmospheric correction, and high-contrast
exoplanet imaging is hard against the bright host star. The JWST offers
far greater sensitivity and stability and a clean point-spread function
from its $6.5\,\si{\meter}$ aperture, but time is scarce and
oversubscribed, the schedule is set by the proposal cycle, and the risk
is the wait and the competition. The recommendation: for bright,
nearby, wide-separation targets the ground telescope now is the
pragmatic choice; for faint or close-in companions where sensitivity and
stability decide, the JWST wait is warranted despite the schedule risk.
\end{exercise}
